
ce q\hphantom{u}i
montre que \(u^{-1}\circ\Gr(\widetilde c)=\Id\).

\textbf{Corollaire 6.12.} \uline{Si} \(K\) \uline{est de caractéristique nulle}, \uline{et} \(\mathord{\wedge}\) \uline{une} \(K\)-\(\lambda\)-\uline{algèbre} \uline{de filtration discrète}, \uline{alors la classe de Chern complétée} :
\[
  \widetilde c:\mathord{\wedge}\longrightarrow \Ch(\Gr\mathord{\wedge})
\]
\uline{est un isomorphisme}.

La démonstration est contenue dans celle de 6.11.

\textbf{Définition 6.13.} \uline{Soit} \(A\) \uline{une} \(K\)-\uline{algèbre graduée}, \uline{et} \(\mathord{\wedge}\) \uline{une} \(K\)-\(\lambda\)-\uline{algèbre} \uline{augmentée}~; \uline{on appelle théorie de Chern de} \(\mathord{\wedge}\) \uline{à valeurs dans} \(A\) \uline{tout} \(\lambda\)-\uline{homomorphisme} : \(\mathord{\wedge}\longrightarrow \Ch(A)\), \uline{compatible avec les augmentations}.

On définit un foncteur sur la catégorie des \(K\)-algèbres graduées, à valeurs dans la catégorie des ensembles, en associant à toute \(K\)-algèbre graduée \(A\) l'ensemble des théories de Chern de \(\mathord{\wedge}\) à valeurs dans \(A\), soit \(\Hom_\lambda(\mathord{\wedge},\Ch(A))\).

\textbf{Proposition 6.14.} \uline{Le foncteur défini en} 6.13 \uline{est représentable}. \uline{Si} \(K\) \uline{est} \uline{de caractéristique nulle}, \uline{et} \(\mathord{\wedge}\) \uline{de filtration discrète}, \uline{il est représenté} \uline{par le gradué associé à} \(\mathord{\wedge}\).

Montrons d'abord le premier point. On considère la \(K\)-algèbre graduée engendrée par des générateurs \(c^i(x)\), avec \(i\ge1\), \(x\in\Fil^1\mathord{\wedge}\), \(c^i(x)\) étant de degré \(i\), et on fait son quotient par les relations suivantes :
\[
\begin{array}{ll}
\text{i)}&
 c^k(x+y)=\displaystyle\sum_{i+j=k}c^i(x)c^j(y),
 \quad\text{pour } k\ge1,\ x,y\in\Fil^1\mathord{\wedge}\ ;\\[0.8em]
\text{ii)}&
 c^k(xy)=P''_k(c^i(x),c^j(y))
 \quad(\text{cf. }6.1),
 \quad\text{pour } k\ge1,\ x,y\in\Fil^1\mathord{\wedge}\ ;\\[0.4em]
\text{iii)}&
 c^k(ax)=R_{k,a}(c^i(x)
 \quad(\text{cf. }6.1),
 \quad\text{pour } k\ge1,\ a\in K,\ x\in\Fil^1\mathord{\wedge}\ ;\\[0.4em]
\text{iv)}&
 c^k(\gamma^j(x))=Q''_{j,k}(c^i(x))
 \quad(\text{cf. }6.1),
 \quad\text{pour } j,k\ge1,\ x\in\Fil^1\mathord{\wedge},
\end{array}
\]
